\documentclass[11pt]{article}
\usepackage[T1]{fontenc}
\usepackage{lmodern}
\usepackage[margin=1in]{geometry}
\usepackage{amsmath,amssymb,amsthm,mathtools}
\usepackage{booktabs,longtable,array}
\usepackage{microtype}
\usepackage{xcolor}
\usepackage{enumitem}
\usepackage{hyperref}
\usepackage{xurl}
\hypersetup{colorlinks=true,linkcolor=blue!45!black,citecolor=blue!45!black,
urlcolor=blue!50!black,pdftitle={Prime separation for Diophantine ZFC consistency},
pdfauthor={Prepared for Vladimir Reshetnikov; AI-assisted research}}
\newtheorem{theorem}{Theorem}[section]
\newtheorem{lemma}[theorem]{Lemma}
\newtheorem{proposition}[theorem]{Proposition}
\newtheorem{corollary}[theorem]{Corollary}
\theoremstyle{definition}
\newtheorem{definition}[theorem]{Definition}
\theoremstyle{remark}
\newtheorem{remark}[theorem]{Remark}
\newcommand{\Z}{\mathbb Z}
\newcommand{\Q}{\mathbb Q}
\newcommand{\N}{\mathbb N}
\newcommand{\ZFC}{\mathsf{ZFC}}
\newcommand{\Con}{\operatorname{Con}}
\newcommand{\RefCode}{\operatorname{Ref}}
\newcommand{\Pos}{\operatorname{Pos}}
\newcommand{\Norm}{\mathcal N}
\newcommand{\vp}{v_3}
\newcommand{\pin}{\nolinkurl{cfbf118b4ad9b2c651fe96b571d9415261766328}}
\newcommand{\code}[1]{\nolinkurl{#1}}
\setlength{\emergencystretch}{2em}
\setlist{itemsep=3pt,topsep=5pt}
\title{Prime separation for Diophantine ZFC consistency\\[5pt]
\large A polynomial in thirteen unrestricted integers\\
of degree 1,392, evaluated in 205 binary operations}
\author{Prepared for Vladimir Reshetnikov and the ProveIt repository\\
\small AI-assisted mathematical research, with exact arithmetic artifacts}
\date{October 23, 2026}

\begin{document}
\maketitle
\begin{abstract}
We construct an explicit ordinary polynomial in thirteen unrestricted
integer unknowns whose lack of integer zeros is provably equivalent in
ZFC to syntactic consistency of ZFC, by composition with the identified
native-compiler and refutation-calibration theorems in ProveIt. Its exact
total degree is \(1,392\), and the supplied shared arithmetic circuit uses
\(205\) binary additions, subtractions, and multiplications. All six fixed
integer coefficients are supplied as exact byte strings.

The new arithmetic method separates several algebraic roots by their
orders at the fixed prime \(3\). A rational sum then forces each root to be
rational, and monicity or an elementary Pell descent recovers integrality.
This permits four Pell and quotient witnesses to share one integer
coordinate. A final step incorporates a root of the positivity guard
itself; a syntactic oddness condition establishes the necessary
nondegeneracy before that guard is recovered. Companion equations have
fourteen, fifteen, and sixteen unrestricted integer unknowns at exact
degrees \(696,244,108\), in \(151,137,132\) operations respectively.

The work retains the source-34 native computational encoding. It does not
claim a globally smallest equation, literature-wide priority, consistency
of ZFC, or a new complete proof-assistant certificate. The new local
reductions, their canonical completeness arguments, their exact degrees,
and their literal polynomial implementations are given in full.
\end{abstract}
\clearpage
\tableofcontents
\clearpage

\section{Results and the meaning of the comparison}
\label{sec:results}

The requested object is a polynomial with fixed integer coefficients and
a specified domain of unknowns. We use syntactic consistency:
\[
 \Con(\ZFC)\quad\equiv\quad
 \neg\exists p\in\N\ \RefCode_{\ZFC}(p),
\]
where \(\RefCode_{\ZFC}\) is the fixed recursive predicate recognizing finite
ZFC refutations used by the repository's calibrated endpoint. A theorem
that ZFC proves an equivalence between this sentence and unsolvability is
different from a proof of either sentence separately.

\begin{theorem}[Main construction]\label{thm:main}
For the six fixed integers specified in Section~\ref{sec:coefficients},
the explicit polynomial \(P_{13}\) of Section~\ref{sec:polynomial}
satisfies
\[
 \ZFC\vdash
 \left(\neg\exists\mathbf x\in\Z^{13}\ P_{13}(\mathbf x)=0\right)
 \ \longleftrightarrow\ \Con(\ZFC).
\]
Its exact total degree is \(1,392\). The delivered arithmetic circuit
evaluates it in \(205\) binary operations. These assertions use the
inherited source contract stated in Section~\ref{sec:contract}.
\end{theorem}

Three companion polynomials retain different numbers of coordinates.
They address different degree and evaluation-cost preferences; the
thirteen-variable equation does not dominate them in every measure.

\begin{table}[ht]
\centering
\begin{tabular}{@{}lrrrr@{}}
\toprule
Polynomial & Integer unknowns & Exact degree & Operations & Final form\\
\midrule
\(P_{13}\) & \textbf{13} & \textbf{1,392} & \textbf{205} & One norm\\
\(P_{14}\) & 14 & 696 & 151 & Norm and guard\\
\(P_{15}\) & 15 & 244 & 137 & Norm and two rows\\
\(P_{16}\) & 16 & 108 & 132 & Four squared rows\\
\bottomrule
\end{tabular}
\caption{All unknowns range over unrestricted integers. Operations count
every live binary \(+,-,\times\), including squaring and multiplication by
fixed numerals. Shared reuse and fixed integer leaves are free.}
\label{tab:main}
\end{table}

The comparison is against the inspected repository snapshot \pin,
including relevant incoming archives, rather than only an earlier
collection README. In the delivered collection, source 67 gives a
thirteen-integer equation of degree \(10,461,744\) in \(519\) operations;
a different earlier thirteen-integer point uses \(518\) operations at a
larger degree~\cite{p67,repo}. The new \(13/1392/205\) point substantially
improves that particular domain and variable-count slice. This is a
bounded comparison to supplied constructions, not a lower-bound theorem
about all polynomials.

The repository also contains equations in ten unrestricted integers and
nine positive integers. Those remain distinct points in the comparison:
fewer variables, the sign restriction on their domains, their degrees,
and their evaluation lengths must all be recorded. We additionally
supply small exact circuit improvements to current ten-integer equations
in Section~\ref{sec:ten}; they preserve the existing polynomials.

\subsection{What changes mathematically}

Previous root mergers needed variable scales, large separating bounds,
or carefully engineered congruences to prevent irrational roots from
canceling. Here a fixed prime already present in the source arithmetic
distinguishes their orders. For four relevant roots the orders are
\[
 2,\quad1,\quad0,\quad-\vp(H).
\]
A fifth guard root has order at least \(3\), or is zero. Every subset
containing an old root therefore has a uniquely smallest order.
A finite sign-product argument shows that such a subset cannot cancel.
A Galois argument converts this fact into rationality of individual
roots.

There are two additional requirements. Rationality must be strengthened
to integrality, and eliminating a guard must not remove a premise needed
to recover that guard. The proof handles the first by monic square
equations and an adjacent-coefficient Pell descent. It handles the second
by making \(c\) odd in the syntax of the polynomial and proving the
necessary bounds before using any positivity guard.

\subsection{What remains inherited}

This is an arithmetic reduction of an established source encoding. The
fixed coefficient tuple still represents the native language associated
with the K215 blank computation. The new thirteen-variable equation
does not invoke the Jones or Sun small-variable finalizer, but its
underlying computational source has not been replaced by a direct
encoding of set-theoretic proof syntax. The source contract below makes
that dependency explicit. No finite numerical test is presented as a
proof of an unbounded computation or of consistency.

\section{The complete thirteen-variable polynomial}
\label{sec:polynomial}

\subsection{Exactly thirteen independent inputs}

Let
\[
 \boxed{j,C,z,h,r,v,w,c_0,T,s,p_2,p_3,x\in\Z.}
 \label{eq:variables}
\]
Every other symbol in this section is a polynomial abbreviation, except
for the single rational abbreviation \(B\), whose denominator is cleared
in the final formula. None is an additional quantified unknown.

Define
\begin{align}
 J&=6j+1,& c&=2c_0+1,& q&=\beta J+1,\nonumber\\
 W&=4(3w+1),& X&=qW,& Y&=(2v+1)q^3,\label{eq:basic}\\
 E&=XY,& a&=(X+1)Y,& H&=4a+3,\nonumber\\
 \Delta&=a^2+H,& L&=(2r+1)c^2,& F_2&=1+\Delta L^2,\nonumber\\
 Q&=\Delta^2L^2,& \delta&=Q-1.&&\nonumber
\end{align}
The native transport and index expressions are
\begin{align}
 F&=(K_0+W)C-z(q-1),\label{eq:transport}\\
 R&=\bigl(q(q-F)-C\bigr)(q^2-1)+(M_C+qM_F)J,\nonumber\\
 k&=R+1+hE,& \eta&=c-Yk,& u&=EYk,\label{eq:rank}\\
 A&=u(u+k)+1,& V&=F_2(cT-R)-c.&&\nonumber
\end{align}
Set
\begin{align}
 D_0&=X+ac,& D_m&=1+\Delta c^2,&
 f_0&=(a+2)L,\label{eq:mainabbrev}\\
 B&=\frac{QV^2-1}{\delta}.&&&&\nonumber
\end{align}
The determinant, interval product, and guard radicand are
\begin{align}
 d&=(a_0F)(b_0C)-(q-a_0F-b_0C)^2,\label{eq:det}\\
 T_\eta&=\eta(k-\eta),\nonumber\\
 \Omega&=dT_\eta-(s^2-d-T_\eta+1)^2
                   -1-p_2^2-p_3(p_3+1).\label{eq:omega}
\end{align}
Here \(a_0,b_0\) are fixed coefficient scales, not unknowns. Finally put
\begin{equation}
 Z=H\bigl[x+18f_0+3(2u+k)+2V\bigr]+2D_0.
 \label{eq:carrier}
\end{equation}

\begin{definition}[Finite polynomial norm]\label{def:norm}
For \(m\ge1\), define
\[
 \Norm_m(Z;A_1,\ldots,A_m)
 =\prod_{\epsilon\in\{-1,1\}^m}
       \left(Z+\sum_{i=1}^m\epsilon_i\sqrt{A_i}\right).
\]
This notation denotes the resulting polynomial in
\(\Z[Z,A_1,\ldots,A_m]\).
\end{definition}

Changing the sign of any formal square root permutes the factors.
Every exponent of that root is therefore even after multiplication.
Consequently Definition~\ref{def:norm} specifies an ordinary polynomial;
it does not require choosing analytic square-root branches.

The promised equation is
\begin{equation}
\boxed{\begin{aligned}
 P_{13}={}&\delta^{16}\Norm_5\bigl(
 Z;\ 324H^2F_2,\ 36H^2A,\ 4H^2B,\ 4D_m,\ 2916H^2\Omega
 \bigr)=0.
\end{aligned}}\label{eq:P13}
\end{equation}
The degree in any one radicand of \(\Norm_5\) is at most \(16\).
Thus \(\delta^{16}\) clears every occurrence of the displayed denominator.
An evaluator using only integer addition, subtraction, and multiplication
is specified in Section~\ref{sec:evaluator}. In particular, the delivered
polynomial has no division operation, variable exponent, unspecified
parameter, or implicit extra witness.

\subsection{The fixed coefficients}
\label{sec:coefficients}

The complete source-34 tuple is used, without mixing its ports with
another native layout:
\[
 \beta=2^{1064921}-1,\qquad
 a_0=2^{1113},\qquad b_0=2^{666104}.
\]
The other three fixed numerals \(K_0,M_C,M_F\) are supplied literally in
the archive. They are the authentic finite native-compiler outputs of
the source-34 layout. \(M_F\) already contains the prescribed shift by
\(\beta\); applying that shift a second time would change the equation.

\begin{table}[ht]
\centering
\begin{tabular}{@{}lrrl@{}}
\toprule
Coefficient & Bits & Raw bytes & Role\\
\midrule
\(\beta\) & 1,064,921 & 133,116 & Native base minus one\\
\(K_0\) & 1,065,043 & 133,131 & Fixed transport coefficient\\
\(M_C\) & 1,064,921 & 133,116 & Native low mask\\
\(M_F\) & 1,064,922 & 133,116 & Already shifted high mask\\
\(a_0\) & 1,114 & 140 & Determinant scale\\
\(b_0\) & 666,105 & 83,264 & Determinant scale\\
\midrule
Total & & 615,883 & 10,607 bytes compressed\\
\bottomrule
\end{tabular}
\caption{Exact unsigned little-endian payloads, compressed individually
with gzip. The binding ledger records full raw and compressed SHA-256 hashes.}
\label{tab:coefficients}
\end{table}

If a decompressed payload consists of bytes \(b_0,\ldots,b_{n-1}\), its
value is exactly \(\sum_{\ell<n}b_\ell256^\ell\). This convention, the
literal payloads, the source pins, and the finite rebuilding recipe bind
all six ports to fixed integers. The table measures the coefficient
inputs to the circuit; it is not the size of a fully expanded polynomial.

Two elementary fixed facts used below are
\[
 \beta>17,\qquad \beta\equiv1\pmod3,\qquad \beta\text{ odd}.
\]
They follow directly from the exponent \(1064921\), which is odd.
No numeral is produced by running a refutation search.

\section{The inherited source contract}
\label{sec:contract}

The new proof has a small arithmetic interface to the existing work.
Stating that interface is necessary: a local polynomial identity alone
does not identify a zero with a ZFC refutation.

\subsection{Soundness of the signed native interface}

Sources 70 and 71 provide, on the same source-34 tuple, the
residue-oriented signed arithmetic system with the calibrated determinant
guard~\cite{p34,p70,p71}. Its variables and quadratic registers can be
projected to the compact interface used here. The retained conditions
include the transport and index equations, first and main Pell norms,
an ordinary strong norm, two auxiliary quotient congruences, the
auxiliary norm, and the sign guards
\[
 d>0,\qquad \eta(k-\eta)>0.
\]
Their theorem states that an unrestricted integer solution yields a
finite syntactic ZFC refutation. The source proof restores its earlier
square guard only after a signed Pell argument recovers the appropriate
power of two; the present work does not reverse that logical order.

For the fifteen- and sixteen-variable companions we use the later
seventeen-integer CRT parent from the incoming Pell-congruences
archive~\cite{incomingpell}. That source proves its quotient compression
before its low-degree projections. The ordinary signed Pell and rank
proofs remain part of that parent's theorem.

\subsection{Canonical completeness}

A refutation supplies canonical source-34 witnesses with
\begin{equation}
 \begin{gathered}
 q=2^t,\quad t\text{ odd},\quad R\ge3,\quad R\equiv3\pmod4,\\
 X=2^R,\quad Y\text{ an odd positive multiple of }q^3,\quad
 c=\psi_{a+2}(R)>2,\\
 \eta>0,\quad k-\eta>0,\quad d>0.
 \end{gathered}\label{eq:canonical}
\end{equation}
Here
\[
 \chi_b(n)+\psi_b(n)\sqrt{b^2-1}
      =(b+\sqrt{b^2-1})^n
\]
defines the standard Pell polynomials for \(b\ge2\).

The source-71 determinant proof applies to the native \(F,C,q\) words
throughout the canonical family, including its permitted rank tuning.
The choice of later strong and auxiliary Pell witnesses does not change
those native words or the interval product \(\eta(k-\eta)\).
This freedom is what permits the normalized completeness construction in
Section~\ref{sec:complete}.

\subsection{Scope of the logical claim}

The existing endpoint connects the K215 blank run with the specified
syntactic refutation predicate. The native compiler connects that run
with the source arithmetic interface. This article proves the new
unbounded implications between that interface and the literal new
polynomials. All arguments use ordinary finite arithmetic, induction,
polynomial algebra, and standard finite extensions of \(\Q\), and can
be formalized in ZFC.

The archive supplies exact finite implementations and mathematical
proofs. It does not export an encoded derivation inside the repository's
object-language ZFC calculus, or transfer the formal status of a different
polynomial to these new ones. Internal independent reviews are identified
as AI-assisted reviews rather than external peer review.

\section{Separation of algebraic roots at the prime three}
\label{sec:separation}

The useful principle concerns rational squares of algebraic numbers.
Write \(\vp(r)\) for the exponent of \(3\) in a nonzero rational number
\(r\), and write \(\Z_{(3)}\) for rationals with denominator prime to \(3\).

\begin{lemma}[No cancellation with a unique least order]\label{lem:nocancel}
Let \(\alpha_1,\ldots,\alpha_n\) be nonzero algebraic numbers with
rational squares. Suppose every subset with at least two members has
a unique member whose squared \(3\)-adic order is smallest, and that
smallest order is even. Then no nonempty signed sum of the
\(\alpha_i\) vanishes.
\end{lemma}
\begin{proof}
A singleton does not vanish. In a larger subset let \(2e\) be the unique
smallest squared order, and divide each root by \(3^e\). Denote the
normalized least root by \(\beta_0\), and the others by
\(\beta_1,\ldots,\beta_s\). Then \(\beta_0^2\) is a unit of
\(\Z_{(3)}\), and each \(\beta_i^2\), \(i>0\), belongs to
\(3\Z_{(3)}\).

Consider the finite product
\[
 U=\prod_{\epsilon\in\{-1,1\}^s}
       \left(\beta_0+\sum_{i=1}^s\epsilon_i\beta_i\right).
\]
Sign invariance in the \(\beta_i\), and simultaneous negation when
changing \(\beta_0\), show that \(U\) is a polynomial with integer
coefficients in all the squared roots. Reduction in \(\Z_{(3)}/3\Z_{(3)}\)
therefore gives
\[
 U\equiv(\beta_0^2)^{2^{s-1}}\not\equiv0\pmod3.
\]
Thus none of the factors can vanish. Every signed sum under consideration
is, up to an overall sign, one of these factors.
\end{proof}

\begin{corollary}[Rational extraction]\label{cor:extract}
Under the hypotheses of Lemma~\ref{lem:nocancel}, if a signed sum of the
roots is rational, each root is rational.
\end{corollary}
\begin{proof}
The field generated by finitely many square roots of rational numbers
is a finite Galois extension of \(\Q\). Each automorphism changes some
root signs. Subtracting its image of the rational sum gives twice a signed
sum of exactly the changed nonzero roots. The lemma excludes such a
nonempty sum. Hence every automorphism fixes every root, and the fixed
field is \(\Q\).
\end{proof}

We will use four squared orders
\begin{equation}
 4,\quad2,\quad0,\quad-2\vp(H),\qquad \vp(H)\ge1,
 \label{eq:orders}
\end{equation}
together with a fifth squared order at least \(6\).
If the fifth root is zero, it is already rational and can be omitted.
If it is nonzero, its order may be odd; this causes no problem because
every subset containing an old root has a unique even minimum, and a
subset consisting only of the fifth root is a singleton.

This proof uses orders of rational numbers and a finite sign product.
It does not need a developed theory of valuations on algebraic
extensions. Nor does it assume that the roots are real: negative
radicands are allowed until rational extraction excludes them at a zero.

\section{Integrality and the joint guard}
\label{sec:integrality}

\begin{lemma}[Rational square integrality]\label{lem:monic}
If \(r\in\Q\) and \(r^2\in\Z\), then \(r\in\Z\).
\end{lemma}
\begin{proof}
Write \(r=a/b\) in lowest terms, \(b>0\). The divisibility
\(b^2\mid a^2\) forces \(b=1\).
\end{proof}

The auxiliary Pell root needs a different argument because its square
is initially a rational number.

\begin{lemma}[Adjacent-coefficient descent]\label{lem:descent}
Let \(Q>1\) and \(V\) be integers. If \(y\in\Q\) satisfies
\[
 QV^2-(Q-1)y^2=1,
\]
then \(y\in\Z\).
\end{lemma}
\begin{proof}
Write \(y=a/m\) in lowest terms with \(m>0\). The equation gives
\(m^2\mid Q-1\). Absolute values preserve it. The case \(V=0\) is
impossible, and \(|V|=1\) gives \(|y|=1\).

For \(V>1,y>0\), define
\[
 V'=(2Q-1)V-2(Q-1)y,\qquad
 y'=(2Q-1)y-2QV.
\]
Since \(m\mid Q-1\), the new \(V'\) is an integer. Expansion preserves
the equation. Moreover \(y>V\), while
\[
 Q^2V^2-(Q-1)^2y^2=QV^2+Q-1>0
\]
gives \(QV>(Q-1)y\). These inequalities imply \(-V<V'<V\); the norm
excludes \(V'=0\). Finally \(2Q-1\equiv1\pmod m\), so \(y'\) has exactly
the same reduced denominator \(m\). Induction on \(|V|\), taking absolute
values at each step, reaches \(|V|=1\) without changing \(m\). Thus \(m=1\).
\end{proof}

This is the adjacent-coefficient descent used in the source-69 and
incoming root-merger constructions~\cite{incomingradical}. The coefficients
\(Q\) and \(Q-1\) are essential; a generic quadratic resultant does not
supply this integrality conclusion.

For the sign conditions, define
\[
 \Pos(p_1,p_2,p_3)=1+p_1^2+p_2^2+p_3(p_3+1).
\]
It is positive for every integer triple and represents every positive
integer. For the latter claim, apply the classical three-square theorem
to \(4n-3\), which is \(1\) modulo \(4\) and is not of the excluded form
\(4^a(8b+7)\). Its three-square representation has one odd square and
two even squares, giving the displayed form after division by four.

\begin{lemma}[Joint positivity]\label{lem:joint}
For integers \(d,T\), the existence of integers \(s,p_1,p_2,p_3\) with
\[
 dT-(s^2-d-T+1)^2=\Pos(p_1,p_2,p_3)
\]
is equivalent to \(d>0\) and \(T>0\).
\end{lemma}
\begin{proof}
The equality gives \(dT>0\), so \(d,T\) have a common nonzero sign.
If \(d=-b,T=-e\) with \(b,e>0\), then
\((s^2+b+e+1)^2>be\), a contradiction.

Conversely let \(d,T>0\) and put \(n=d+T-1\ge1\). Choose an integer
square nearest \(n\). Its distance from \(n\) has square strictly less
than \(n\): between \(m^2\) and \((m+1)^2\), the nearest distance is at
most \(m\), and the only nonzero-distance cases have \(n>m^2\).
Since \(n\le dT\), the residual \(dT-(s^2-n)^2\) is a positive integer.
Represent it by \(\Pos\).
\end{proof}

In our notation, recovering an integer \(p_1\) with
\(p_1^2=\Omega\) restores the joint guard
\begin{equation}
 G=dT_\eta-(s^2-d-T_\eta+1)^2-\Pos(p_1,p_2,p_3)=0.
 \label{eq:G}
\end{equation}
The important feature is that the new equation will recover this guard
after the other roots, using premises that do not depend on \(G\).

\section{Bounds and residues before using the guard}
\label{sec:unconditional}

\begin{lemma}[Unconditional nondegeneracy]\label{lem:bounds}
For every assignment of the thirteen integer inputs in
Section~\ref{sec:polynomial},
\[
 H\ne0,\qquad \Delta>0,\qquad L\ne0,\qquad Q>1,\qquad \delta\ne0,
 \qquad A\ge1.
\]
Furthermore,
\[
 a\equiv H\equiv\Delta\equiv0\pmod3,\qquad
 F_2\equiv A\equiv D_m\equiv1\pmod3,\qquad Q\equiv0\pmod9,
\]
and \(B\) is a unit congruent to \(1\) modulo \(3\) in \(\Z_{(3)}\).
\end{lemma}
\begin{proof}
The odd integer \(J=6j+1\) is nonzero, so \(q=\beta J+1\) is even and
\(|q|\ge\beta-1\ge16\). The nonzero multiple of four
\(W=4(3w+1)\) has \(|W|\ge4\). Hence
\[
 |X|\ge64,\quad |Y|\ge4096,\quad
 |a|=|X+1|\,|Y|\ge63\cdot4096.
\]
Thus \(\Delta=(a+1)(a+3)>0\). Because \(c\) and \(2r+1\) are odd,
\(L=(2r+1)c^2\ne0\), and \(Q=\Delta^2L^2>1\).
The integer \(H=4a+3\) cannot be zero and is odd. Also \(Y\) is even,
so \(a\) and hence \(\Delta-1\) are even.

Modulo \(3\), the fixed coefficient and affine restrictions give
\[
 J\equiv1,\quad q\equiv2,\quad W\equiv1,\quad X\equiv2,\quad a\equiv0.
\]
The claimed congruences for \(H,\Delta,F_2,D_m,Q\) follow.
Writing \(b=XY^2\), we have
\[
 u=bk,\qquad A=b(b+1)k^2+1.
\]
For integer \(b\), \(b(b+1)\ge0\). Modulo \(3\), \(b\) is either \(0\)
or \(2\), so \(b(b+1)\equiv0\). This proves the assertions about \(A\).
Finally both \(QV^2-1\) and \(Q-1\) are \(-1\) modulo \(3\), so their
ratio \(B\) is the stated unit.
\end{proof}

The proof neither uses the first or auxiliary Pell equation nor assumes
the joint guard. In particular \(c=\pm1\) is permitted at this stage:
nonzero oddness is sufficient for the preliminary bounds. Such values
will be excluded, if necessary, when the recovered interval guard is
combined with the large value of \(|Y|\).

\section{Soundness of the thirteen-variable equation}
\label{sec:soundness}

Suppose \(P_{13}=0\). By Lemma~\ref{lem:bounds}, the denominator-clearing
factors \(H,\delta\) are nonzero. The finite norm product therefore
supplies algebraic numbers \(f,\tau,y,D,p_1\), with signs absorbed into
their names, satisfying
\begin{equation}
 f^2=F_2,\quad \tau^2=A,\quad y^2=B,\quad
 D^2=D_m,\quad p_1^2=\Omega,
 \label{eq:recovered}
\end{equation}
and
\begin{equation}
 \frac{Z}{2H}=9f+3\tau+y+\frac{D}{H}+27p_1.
 \label{eq:rationalsum}
\end{equation}
Repeated roots, reducible quadratics, and \(\Omega=0\) are allowed in this
factorization.

The squared orders of the first four terms in
\eqref{eq:rationalsum} are precisely those in \eqref{eq:orders}.
The fifth squared order is at least \(6\), unless \(p_1=0\).
Section~\ref{sec:separation} makes every root in
\eqref{eq:recovered} rational. Lemma~\ref{lem:monic} makes
\(f,\tau,D,p_1\) integers; Lemma~\ref{lem:descent} makes \(y\) an integer.
This also shows that any initially negative radicand is impossible at a
zero. It is unnecessary, and would be circular, to assume its positivity
before root recovery.

Set
\[
 \gamma=\frac{D-D_0}{H}.
\]
Substitute the definition of \(Z\) into \eqref{eq:rationalsum}. All terms
except the displayed factor \(1/2\) are integers, so \(2\gamma\in\Z\).
Also \(H\gamma=D-D_0\in\Z\). Since \(H\) is odd,
\(\gcd(2,H)=1\) and \(\gamma\in\Z\). The midpoint carrier thus creates no
unpaid parity condition.

The integer \(p_1\) restores \(G=0\), and
Lemma~\ref{lem:joint} gives \(d>0\) and
\(\eta(k-\eta)>0\). The source arithmetic variables can now be restored
explicitly.

\subsection{Polynomial restoration of the quotients}

Put \(i=2r+1\). The ordinary strong ordinate and scale are
\[
 I_{\rm old}=\Delta L,\qquad
 i_{\rm old}=\Delta i,\qquad
 r_{\rm old}=\frac{i_{\rm old}-1}{2}.
\]
Both \(\Delta\) and \(i\) are odd, so \(r_{\rm old}\) is an integer.
Moreover
\[
 I_{\rm old}=i_{\rm old}c^2,\qquad
 I_{\rm old}^2=Q=\Delta(f^2-1).
\]
The normalized source makes the two quotient restorations especially
simple:
\begin{equation}
 o=f(cT-R),\qquad
 j_{\rm aux}=F_2T-1-\Delta i^2c^3R.
 \label{eq:quotientrestoration}
\end{equation}
Direct multiplication gives
\[
 V+c=fo,\qquad V+R=cj_{\rm aux}.
\]
There is no unproved division by \(c\): the identity
\(F_2-1=\Delta i^2c^4\) has already supplied the required quotient as a
polynomial.

\subsection{Restoration of the remaining source rows}

Restore
\[
 \alpha=q-F-C,\quad \zeta=k-\eta,\quad P=k+u,\quad D_{\rm old}=D.
\]
Then
\[
 \tau^2-P^2+Pk-1=\tau^2-u(u+k)-1=0,
 \qquad D_{\rm old}^2-\Delta c^2-1=0.
\]
The native transport row follows from \eqref{eq:transport}. The index
row follows from \(k=R+1+hE\). The source's elementary product registers
are restored by their defining integer polynomials:
\[
 q^2-1,\quad q^3,\quad c^2,\quad f^2,\quad Q,\quad y^2-V^2.
\]
The auxiliary equations follow from
\(QV^2-(Q-1)y^2=1\) and \eqref{eq:quotientrestoration}.
For source 70's affine coordinate use
\(J_{3,\rm old}=(J-1)/3=2j\); the \(w\) coordinate is unchanged.
Finally represent the recovered positive guard quantities by their
private \(\Pos\) triples.

Every restored coordinate is an integer, and the source's original rows
vanish. The signed native soundness theorem therefore supplies a finite
ZFC refutation. This establishes the forward existential implication
without assuming either consistency or the absence of a refutation.

\section{Canonical completeness and the normalized strong root}
\label{sec:complete}

Begin with a finite ZFC refutation and take the canonical family
\eqref{eq:canonical}. Its native
\(J=(q-1)/\beta\) is a finite repunit in the even base
\(2^{1064921}\), hence is odd. Residue orientation gives
\(J\equiv1\pmod3\), so \(J=6j+1\) for an integer \(j\).
The main Pell parameter \(a+2\) is even. The Pell recurrence modulo \(2\)
then shows that \(c=\psi_{a+2}(R)\) is odd because \(R\) is odd.
Thus \(c=2c_0+1\) is a permissible syntactic restriction on this family.

We must prove that the stronger normalization \(L=(2r+1)c^2\) does not
destroy completeness.

\begin{lemma}[Odd normalized strong choice]\label{lem:normalized}
Under \eqref{eq:canonical}, there are positive integers \(f,L\) with
\[
 f^2-\Delta L^2=1,\qquad L/c^2\text{ an odd integer}.
\]
\end{lemma}
\begin{proof}
Let \(b=a+2\), choose \(m=cR\), and set
\(f=\chi_b(m)\), \(L=\psi_b(m)\).
The Pell norm gives the first equation. Composition gives
\[
 L=\psi_b(R)\psi_{\chi_b(R)}(c)=c\psi_D(c),
 \qquad D=\chi_b(R).
\]
Since \(D^2-\Delta c^2=1\), we have \(D^2\equiv1\pmod c\).
The recurrence consequently gives
\[
 \psi_D(n)\equiv nD^{n-1}\pmod c.
\]
At \(n=c\) this proves \(c^2\mid L\). Both \(c\) and \(R\) are odd, so
\(m\) is odd. As \(b\) is even, the recurrence modulo \(2\) makes \(L\)
odd. Its integer quotient by the odd square \(c^2\) is therefore odd.
\end{proof}

Choose \(r\) with \(L/c^2=2r+1\). The remaining issue is the stronger
form of the auxiliary quotient, which contains \(f^2\) rather than \(f\).

\begin{lemma}[Fresh auxiliary witnesses]\label{lem:auxcomplete}
Let \(R\ge3\), \(R\equiv3\pmod4\), \(c=\psi_b(R)\), and suppose
\(f,L>0\) satisfy \(f^2-(b^2-1)L^2=1\) with \(c\mid L\).
Put \(S=(b^2-1)L\) and \(Q=S^2\). There are positive integers \(T,y\)
such that
\[
 V=f^2(cT-R)-c,\qquad QV^2-(Q-1)y^2=1.
\]
\end{lemma}
\begin{proof}
This is source 61's positive auxiliary reconstruction~\cite{p61}.
Its explicit choice is
\[
 V_*=\frac{\chi_S(R)}{S},\qquad y_*=\psi_S(R).
\]
The odd-index Pell identities make \(V_*\) an integer and give
\[
 V_*\equiv-c\pmod{f^2},\qquad V_*\equiv-R\pmod c.
\]
Here \(S^2\equiv1-b^2\pmod{f^2}\), and
\(R\equiv3\pmod4\) fixes the required signs.
The strong equation gives \(f^2\equiv1\pmod c\), hence
\(\gcd(f^2,c)=1\). Therefore
\[
 T=\frac{V_*+c+Rf^2}{cf^2}
\]
is an integer and has \(f^2(cT-R)-c=V_*\).
The source identities give the auxiliary norm with \(y=y_*\);
the chosen positive branches make \(T,y\) positive.
\end{proof}

The two lemmas choose fresh strong and auxiliary witnesses. They do not
claim that every earlier source solution had the stronger quotient
divisibility. The new choices leave \(q,F,C,k,\eta\) unchanged. Thus the
canonical determinant \(d>0\) and interval product
\(T_\eta>0\) remain valid. Lemma~\ref{lem:joint} supplies
\(s,p_1,p_2,p_3\).

Let \(\tau,\gamma\) be the canonical first root and main congruence
quotient. Define the integer
\begin{equation}
 x=18f+6\tau+2y+2\gamma+54p_1
       -18(a+2)L-3(2u+k)-2V.
 \label{eq:completecarrier}
\end{equation}
Substitution makes one factor in \eqref{eq:P13} vanish. Therefore a finite
refutation supplies an integer zero of \(P_{13}\). Together with
Section~\ref{sec:soundness}, this proves Theorem~\ref{thm:main}.

\section{Three companion constructions}
\label{sec:companions}

\subsection{Fourteen integers}

Restore \(p_1\) as an independent coordinate, retain \(c\) itself as an
independent integer, and omit the syntactic substitution \(c=2c_0+1\).
Use the same normalized strong expressions and the guard \(G\) in
\eqref{eq:G}. Define the raw four-root polynomial
\[
 N_{14}=\delta^8\Norm_4
 \bigl(Z;324H^2F_2,36H^2A,4H^2B,4D_m\bigr).
\]
Then set
\[
 P_{14}=N_{14}+G^2+N_{14}G^2.
\]
The equation is equivalent to \(N_{14}=G=0\), since
\((1+N_{14})(1+G^2)=1\) and the second factor is a positive integer.
Here the guard is available before root recovery. It gives
\[
 |k|\ge2,\quad0<|\eta|<|k|,\quad
 |c|=|Yk+\eta|>(|Y|-1)|k|>1.
\]
Thus \(L\ne0,Q>1\), and the preceding four-root argument applies.
Its exact degree is \(696\), and its compact circuit has \(151\) operations.

\subsection{Sixteen integers and the fixed multiplier three}

For the remaining companions retain an independent integer \(f\) and put
\[
 J=3j+1,\quad L=(2r+1)c^2,\quad Q=L^2,\quad F_2=f^2,
 \qquad V=fcT-f^2R-c.
\]
Keep the separate strong row
\(S_0=Q-\Delta(f^2-1)\).
All other common abbreviations have the same meanings. The
seventeen-integer CRT parent supplies this interface~\cite{incomingpell}.

The guard gives \(|c|>1,Q>1\). Since \(\Delta\equiv0\pmod3\), the strong
row makes \(3\mid L\) and \(9\mid Q\). Thus the fixed multiplier \(3\)
separates the first and auxiliary roots. Put
\begin{align*}
 t_0&=2u+k,& \ell_0&=k^2-4,\\
 J_*(x)&=\delta x^2+\delta(6t_0+4V)x
                         +9\delta\ell_0-4(V^2-1),\\
 H_*&=12\delta t_0V,\\
 N_{16}(x)&=J_*(x)^2+2H_*J_*(x)
       +576\delta A(1-V^2)+144\delta^2\ell_0V^2.
\end{align*}
An exact identity gives
\[
 N_{16}(x)=\delta^2\Norm_2(x+3t_0+2V;36A,4B).
\]
This replaces \(\tau,y\) by \(x=6\tau+2y-3t_0-2V\).
Root separation, monicity, and Lemma~\ref{lem:descent} prove the converse.
The sum of squares of \(N_{16}\), \(S_0\), \(G\), and
\((D_0+H\gamma)^2-D_m\) is \(P_{16}\), of degree \(108\) in
sixteen integers and \(132\) operations.

\subsection{Fifteen integers}

Eliminate \(\gamma\) as well, using
\[
 N_{15}(x)=H^8
 \prod_{(H\gamma+D_0)^2=D_m}N_{16}(x-2\gamma).
\]
The product includes both roots and denotes its denominator-free
polynomial expansion. The added weighted root \(D/H\) has order
\(-\vp(H)\), below the first two. The same rational extraction and
odd-\(H\) midpoint argument recover an integer \(\gamma\).
With \(S=S_0^2+G^2\), set
\[
 P_{15}=N_{15}+S+N_{15}S.
\]
This equation is equivalent to the three required rows. Its exact degree
is \(244\), and its circuit has \(137\) operations. These two companions
retain the existing CRT parent's rank and divisibility proof; those
dependencies are not replaced by finite samples.

\section{Division-free evaluation and arithmetic counts}
\label{sec:evaluator}

A norm product is a compact algebraic specification, but a concrete
straight-line evaluator makes the equation independently reproducible.
Here is the four-root evaluator used inside the thirteen-variable circuit.

For inputs \(Z,A_1,A_2,A_3,C_*,\delta\), set
\begin{align*}
 \sigma&=A_1+A_2+A_3,&
 \pi&=A_1A_2+(A_1+A_2)A_3,&
 \rho&=A_1A_2A_3,\\
 a_*&=\delta(Z^2-\sigma)+C_*,&
 t_*&=\delta C_*Z^2,\\
 k_*&=a_*^2+4(t_*-\delta^2\pi),&
 r_*&=16\delta^3\rho,\\
 b_*&=a_*k_*-r_*,\\
 e_*&=k_*^2+16t_*a_*^2-4r_*(\delta Z^2+C_*),\\
 \mathcal Q&=e_*^2-64t_*b_*^2.
\end{align*}
Direct multiplication of formal norm factors proves
\[
 \mathcal Q
   =\delta^8\Norm_4(Z;A_1,A_2,A_3,C_*/\delta).
\]
Use
\[
 A_1=324H^2F_2,\quad A_2=36H^2A,\quad A_3=4D_m,\quad
 C_*=4H^2(QV^2-1).
\]
Evaluate \(\mathcal Q\) at \(Z-\omega\) in the quadratic algebra
\(\omega^2=2916H^2\Omega\). If its result is \(E_*+O_*\omega\), then
\[
 P_{13}=E_*^2-2916H^2\Omega\,O_*^2.
\]
For this calculation addition is componentwise and multiplication is
\[
 (e,o)(e',o')=(ee'+2916H^2\Omega\,oo',\ eo'+oe').
\]
The pair entries are derived polynomial expressions; they do not add
unknowns. This gives an entirely division-free definition of
\eqref{eq:P13}.

The implementation shares repeated nodes, uses fixed powers only, and
expands every power into multiplication. The compact circuits have the
following verified operation splits:
\begin{center}
\begin{tabular}{@{}lrrrr@{}}
\toprule
 & Additions & Subtractions & Multiplications & Total\\
\midrule
\(P_{13}\) & 47 & 28 & 130 & 205\\
\(P_{14}\) & 38 & 23 & 90 & 151\\
\(P_{15}\) & 33 & 24 & 80 & 137\\
\(P_{16}\) & 33 & 24 & 75 & 132\\
\bottomrule
\end{tabular}
\end{center}
These are exact counts of supplied evaluators, not proofs of minimum
circuit size or estimates of bit complexity. Operations on million-bit
fixed integers are not unit-cost in an implementation concerned with
running time. A separate coefficient-based implementation expands the
successive quadratic norms in their carrier coordinate and provides an
independent route for verification and degree certification.

\section{Exact degree proofs}
\label{sec:degrees}

Naive degree propagation through a compact circuit can overestimate the
degree because it misses cancellations. The compact fourteen-variable
circuit, for example, has a propagated bound of \(792\), while its actual
degree is \(696\). We use the independent coefficient construction to
prove upper bounds, and exact finite-field specializations to certify
nonzero leading coefficients after fixing the actual numerical ports.

\subsection{Sixteen and fifteen variables}

The coefficient construction gives upper bounds \(54\) for the norm
\(N_{16}\) and \(156\) for \(N_{15}\). The former is the largest residual
degree in its sum of squares, giving degree at most \(108\). In the
fifteen-variable finalizer, \(G\) has degree \(44\), so
\(S=S_0^2+G^2\) has degree \(88\), and \(N_{15}+S+N_{15}S\) has degree at
most \(244\). The matching nonzero certificates below make these bounds
exact.

\subsection{The centered strong root and degree 696}

Immediately before the final strong-root norm in the fourteen-variable
coefficient construction, write the degree-eight carrier polynomial as
\(\sum_{i=0}^8a_i x^i\). Direct arithmetic propagation bounds the
coefficient degrees by
\[
 (292,295,294,277,260,237,214,191,168).
 \label{eq:coefficient8}
\]
The final root is centered at \(f_0=(a+2)L\). Writing \(w_f=f-f_0\),
its exact equation is
\begin{equation}
 w_f^2+2f_0w_f+(L^2-1)=0,
 \qquad \deg f_0\le9,\quad \deg(L^2-1)\le6.
 \label{eq:centeredstrong}
\end{equation}
The identity \(f_0^2-F_2=L^2-1\) is responsible for the degree saving.

Scale every independent coordinate along a fixed real direction by a
parameter \(t\to+\infty\). On a nonempty open set of directions, the
degree-nine leading coefficient of \(f_0\) is nonzero. The roots of
\eqref{eq:centeredstrong} are eventually real: \(f_0^2\) has degree
eighteen while \(L^2-1\) has degree six. One root is \(O(t^9)\). By the
product of the roots, or the conjugate quotient formula, the other is
\(O(t^{-3})\). The two arguments \(x-18w_f\) are therefore
\(O(t^9)\) and \(O(t)\). Their polynomial values have product of order
at most
\[
 t^{M_9+M_1},\qquad
 M_9=\max_i(\deg a_i+9i)=312,\quad
 M_1=\max_i(\deg a_i+i)=296.
\]
Thus the raw norm has degree at most \(608\). This is a global degree
bound: a homogeneous polynomial component of larger degree would vanish
on the open set of directions just used, and hence would be identically
zero. The argument does not merely estimate the polynomial at one
chosen numerical ray.

The guard has exact degree \(44\), since its leading term is the negative
square of the degree-22 product \(\eta(k-\eta)\). Therefore
\[
 \deg P_{14}\le608+2\cdot44=696.
\]

\subsection{The guard norm and degree 1,392}

Before folding the guard root, the degree-sixteen raw norm as a
polynomial in its carrier has the following coefficient-degree bounds:
\[
\begin{split}
 (&612,607,606,597,588,571,554,537,520,\\
  &497,474,451,428,405,382,359,336).
\end{split}
\]
The affine substitution \(c=2c_0+1\) does not increase them.
Since \(\deg\Omega\le44\), assign weight \(22\) to a formal root
\(\omega^2=\Omega\). Either substituted factor has weighted degree at
most
\[
 \max_{0\le i\le16}\bigl(\deg a_i+22i\bigr)=696.
\]
Its quadratic norm therefore has degree at most \(1,392\).
This can also be read as direct degree propagation in the quadratic
algebra; it does not require that the formal root already be real or
integral.

\subsection{Nonzero leading coefficients on the actual coefficient slice}

The checker replaces each integer coordinate by a recorded integer slope
times \(t\), reduces all six actual coefficient payloads modulo the prime
\(65521\), and performs exact polynomial convolution. A nonzero
coefficient at the upper-bound degree certifies that the corresponding
homogeneous component is not zero.

\begin{table}[ht]
\centering
\begin{tabular}{@{}lrrrr@{}}
\toprule
 & Norm degree & Norm coefficient & Final degree & Final coefficient\\
\midrule
16 integers & 54 & 17,367 & 108 & 19,526\\
15 integers & 156 & 29,172 & 244 & 16,359\\
14 integers & 608 & 32,234 & 696 & 50,632\\
13 integers & 1,392 & 56,423 & 1,392 & 56,423\\
\bottomrule
\end{tabular}
\caption{Nonzero leading coefficients modulo \(65521\). Each row uses its
own fully recorded slope specialization.}
\label{tab:degreecert}
\end{table}

The prime is checked by trial division. The convolution uses arbitrary
precision integers before modular reduction. The exact slope vectors,
coefficient residues, intermediate coefficient bounds, and outputs are
included. These are rigorous finite lower certificates, not probabilistic
claims of polynomial identity. Together with the symbolic upper bounds
they establish all exact degrees in Table~\ref{tab:main}.

\section{Verification and reproducibility}
\label{sec:verification}

The artifact has two deliberately different implementations of the norm
construction. The compact one uses the closed four-root formula and
quadratic pair arithmetic. The coefficient implementation treats the
carrier as a polynomial variable and takes successive quadratic norms.
The identities defining both are proved algebraically. Independent
formal sign-product expansions check the universal norm formulas, while
modular comparisons with arbitrary coordinate values check transcription
of the literal circuits.

The additional checker
\code{independent_audit/check_primary_formula.py} imports neither circuit
builder. It independently states all twenty-eight source abbreviations
and the closed norm evaluator, then canonicalizes the actual 205-gate
packet by elementary ring identities. This proves exact equality with
the displayed \(P_{13}\), without sampling. In particular it checks the
passage from \(\delta^8\) to \(\delta^{16}\), all five radicand scales,
the guard fold, and agreement of the actual source-34 coefficient bytes.

The finite checks establish coefficient integrity, declared variable
counts, absence of undeclared inputs, node liveness, operation counts,
exact algebraic identities, and degree certificates. Normal and optimized
Python runs retain the active checks. They do not search for a solution
of the consistency equation or prove its unsolvability.

The archive includes the complete LaTeX article, its compiled PDF,
the two circuit implementations, literal polynomial packets, coefficient
payloads and rebuilding data, provenance and license notices, verification
scripts, and receipts. Its source snapshot is fixed even if the repository
subsequently advances. The ordinary proof in this article is also retained
in a readable proof file so that the algebraic construction can be reviewed
independently of the PDF layout.

A rebuild should run the included verification driver before rebuilding
the report. From the extracted root, use
\texttt{python3 -O verify\_all.py}; its checks remain active under Python
optimization. The mathematical checks use ordinary Python; the report uses
a standard \texttt{pdflatex} installation and two passes. PDF byte identity
is not required, since creation metadata and TeX versions can vary.
Polynomial packet identity and coefficient hashes are the stable
reproducibility targets.

\section{Ancillary exact improvements to ten-integer evaluators}
\label{sec:ten}

A separate part of the work applies three exact identities to the
current ten-integer shear circuits~\cite{shears,p68}. These changes preserve
each polynomial identically, so its witness domain, consistency theorem,
and exact degree are inherited without a new semantic reduction.

The first reuses already available powers:
\[
 S^5=(S^2)(S^3).
\]
The second changes the placement of a two-root norm:
\begin{align*}
 &(Z^2-\Sigma^2(a_0+a_1))^2-4\Sigma^4a_0a_1\\
 &\hspace{10mm}=(Z^2+\Sigma^2(a_0-a_1))^2
                         -4Z^2\Sigma^2a_0.
\end{align*}
The third factors a common block; writing \(c_*\) for the existing
capacity register and \(D=C-KL\), it is
\begin{align*}
 &8c_*^3gK^2-g^2\bigl(32D^2c_*^3+g^2K^2\bigr)\\
 &\hspace{10mm}=8gc_*^3(K^2-4gD^2)-g^4K^2.
\end{align*}
Each identity saves one operation in the supplied shared evaluator,
for three operations per circuit.

\begin{center}
\begin{tabular}{@{}lrrr@{}}
\toprule
Variant & Exact degree & Prior operations & New operations\\
\midrule
Predecessor & 1,245,324 & 480 & 477\\
Linear shear & 1,179,788 & 484 & 481\\
Cubic shear & 1,114,252 & 487 & 484\\
Six-term shear & 1,081,125 & 521 & 518\\
\bottomrule
\end{tabular}
\end{center}

The optimization packets retain source hashes and the exact substituted
circuits. Modular checks are used only as transcription diagnostics;
the three displayed ring identities are the proofs. These small
improvements are separate from the thirteen-variable prime-separation
construction and should not be described as the same type of advance.

\section{Further research questions}
\label{sec:future}

\paragraph{How far can local-order separation be generalized?}
The proof uses the prime \(3\) because the source's existing affine
restrictions make \(a,H,\Delta\) divisible by \(3\), while the relevant
radicands are units. Other encodings may expose a different prime or
several primes. A useful general result would minimize fixed root weights
subject to a proved noncancellation condition, while accounting for
the operation cost of those weights and the cost of new congruence
restrictions.

\paragraph{Can a second guard coordinate be removed?}
The fifth root works because one guard coordinate appears as a separate
square and the other arithmetic roots remain recoverable without the
guard. Repeating that argument is not automatic: the remaining guard
coordinates share a single quadratic representation, so they do not
supply independent radicands. A further reduction needs a new complete
parameterization or another arithmetic invariant, rather than simply
adding more formal norm factors.

\paragraph{Can normalized strong coordinates help other degree slices?}
The choice \(L=(2r+1)c^2\) makes quotient restoration polynomial and
removes a preliminary rank dependency from the new soundness proof.
It also changes the degrees of \(F_2,Q,V\). Retaining one intermediate
coordinate might lower the degree or operation count at fifteen or
sixteen variables. Such variants should be compared as complete literal
circuits, including every restored divisibility or sign condition.

\paragraph{Can degree certificates be extracted automatically from norms?}
The centered strong-root argument is a small example of using the
different growth rates of conjugate roots to prove cancellation-aware
degree bounds. An exact tool that combines coefficient bounds, Newton
polygons, and a finite nonzero specialization could make such bounds
routine. Its output should distinguish a symbolic upper-bound proof
from a numerical lower certificate, as the present construction does.

\paragraph{What is the appropriate notion of simplicity?}
Variable count, degree, circuit operations, expanded monomial count,
coefficient height, and total encoded description length are different
objectives. The current native coefficient inputs occupy 615,883 raw
bytes, although they compress well. A comparison in a fixed description
language should count the coefficient recipe or the literal bytes
consistently. It should not infer bit complexity from the number of
arithmetic gates.

\paragraph{Can the computational source be replaced?}
The user explicitly allowed non-K215 and non-Jones routes. The present
advance is in the arithmetic representation of a calibrated source.
A direct proof-language or equiconsistent-theory encoding would need
a complete source predicate, a proved bridge to the same syntactic
consistency sentence, and all of its sequence and arithmetic encodings
materialized before a fair size comparison. The prime-separation lemma
could be useful in such a compiler without forcing it to reuse this
native tableau.

\paragraph{What should be formalized next?}
The elementary sign-product lemma, rational extraction, adjacent Pell
descent, and midpoint integrality are focused independent targets.
The literal norm identities and quotient restorations are suitable for
a polynomial reflection checker. A complete formal result would then
need to connect those statements to the actual coefficient bytes and
to the inherited native and syntactic-calibration theorems. Closing that
boundary is a separate deliverable from the finite circuit audits
reported here.

\begin{thebibliography}{12}
\raggedright

\bibitem{repo}
V.~Reshetnikov and ProveIt contributors.
\emph{ProveIt: DiophantineConsistency, HilbertTenthProblem, and incoming
research packages}. Inspected snapshot
\pin, October 23, 2026.
\url{https://github.com/VladimirReshetnikov/ProveIt/tree/cfbf118b4ad9b2c651fe96b571d9415261766328}.

\bibitem{p34}
ProveIt contributors.
\emph{Source 34: integer-variable consistency forms and the
mixed-parity-unrounded-1064921 native tuple}.
\url{https://github.com/VladimirReshetnikov/ProveIt/tree/cfbf118b4ad9b2c651fe96b571d9415261766328/Computability/DiophantineConsistency/con-zfc-diophantine-forms/34-integer14}.
The whole tuple and its finite rebuilding recipe are the coefficient
source used here.

\bibitem{p61}
ProveIt contributors.
\emph{Source 61: congruence carriers}.
In particular, the direct normalized rank and positive auxiliary
reconstruction lemmas.
\url{https://github.com/VladimirReshetnikov/ProveIt/tree/cfbf118b4ad9b2c651fe96b571d9415261766328/Computability/DiophantineConsistency/con-zfc-diophantine-forms/61-congruence-carriers}.

\bibitem{p67}
ProveIt contributors.
\emph{Source 67: oriented guards}.
Thirteen-integer comparison forms and signed arithmetic dependencies.
\url{https://github.com/VladimirReshetnikov/ProveIt/tree/cfbf118b4ad9b2c651fe96b571d9415261766328/Computability/DiophantineConsistency/con-zfc-diophantine-forms/67-oriented-guards}.

\bibitem{p68}
ProveIt contributors.
\emph{Source 68: joint positivity}.
Ten-integer evaluator and coefficient conventions inherited by the
ancillary optimization branch.
\url{https://github.com/VladimirReshetnikov/ProveIt/tree/cfbf118b4ad9b2c651fe96b571d9415261766328/Computability/DiophantineConsistency/con-zfc-diophantine-forms/68-joint-positivity}.

\bibitem{p70}
ProveIt contributors.
\emph{Source 70: residue orientation}.
Signed power recovery, canonical residue restrictions, and exact
integer-coordinate projections.
\url{https://github.com/VladimirReshetnikov/ProveIt/tree/cfbf118b4ad9b2c651fe96b571d9415261766328/Computability/DiophantineConsistency/con-zfc-diophantine-forms/70-residue-orientation}.

\bibitem{p71}
ProveIt contributors.
\emph{Source 71: calibrated determinant guards}.
The common-sign determinant and its canonical completeness on source 34.
\url{https://github.com/VladimirReshetnikov/ProveIt/tree/cfbf118b4ad9b2c651fe96b571d9415261766328/Computability/DiophantineConsistency/con-zfc-diophantine-forms/71-determinant-guards}.

\bibitem{incomingradical}
ProveIt incoming research.
\emph{Scaled-root compression for Diophantine consistency},
archive \code{con-zfc-radical-residue-2026-10-23}.
The integer proof, circuit utility, and six exact coefficient payloads
are retained with their source identities in the present archive.
Inspected under \code{docs/incoming} at the snapshot in~\cite{repo}.

\bibitem{incomingpell}
ProveIt incoming research.
\emph{Pell congruences for Diophantine consistency},
archive \code{con_zfc_pell_congruences_2026-10-23}.
The CRT quotient proof and \code{transport/degree88_17.json} provide the
seventeen-integer parent for the companions.
Inspected under \code{docs/incoming} at the snapshot in~\cite{repo}.

\bibitem{shears}
ProveIt incoming research.
\emph{Polynomial shears for Diophantine consistency},
archive \code{con_zfc_polynomial_shears}.
The four ten-integer circuits in its \code{norm10} component are the
parents of the exact three-operation rewrites.
Inspected under \code{docs/incoming} at the snapshot in~\cite{repo}.

\end{thebibliography}
\end{document}

